\documentclass[10pt,a4paper]{amsart}
\usepackage[margin=19mm]{geometry}
\usepackage{amsmath,amssymb,mathtools}
\usepackage[scr=rsfs,cal=euler]{mathalfa}
\usepackage{xfrac}
\usepackage[unicode,hidelinks,bookmarks=false]{hyperref}
\newcommand{\ie}{i.e.\ }
\newcommand{\cf}{cf.\ }
\newcommand{\resp}{respectively\ }
\newcommand{\Subsection}{\subsection}
\providecommand{\nobreakdash}{\nobreak\hbox{-}}
\setlength{\emergencystretch}{2em}
\multlinegap0pt
\usepackage[shortlabels]{enumitem}
\usepackage{amssymb}
\usepackage{mathtools}
\usepackage{xcolor}
\usepackage{float}
\usepackage{tikz}
\newtheorem{corollary}{Corollary}
\newtheorem{lemma}{Lemma}
\newtheorem{proposition}{Proposition}
\DeclareMathOperator{\Ind}{Ind}
\newcommand{\K}{\mathbb K}
\DeclareMathOperator{\del}{del}
\DeclareMathOperator{\link}{link}
\title{Vertex Dismissibility and Scalability of Simplicial Complexes}
\author{Mohammed Rafiq Namiq}
\date{Source: arXiv:2603.10736v3 (CC BY 4.0)}
\begin{document}
\maketitle

\noindent\emph{Source and licence.} Vertex Dismissibility and Scalability of Simplicial Complexes. Version arXiv:2603.10736v3 (CC BY 4.0), distributed by arXiv under the Creative Commons Attribution 4.0 International licence, http://creativecommons.org/licenses/by/4.0/, as recorded in the arXiv licence metadata for this exact version and shown on the abstract page https://arxiv.org/abs/2603.10736v3. Full licence notice: \texttt{licenses/arxiv-2603.10736-CC-BY-4.0.txt}.\par\smallskip

\noindent\textit{Editorial scope.} Counted unit: the article's proposition prop:7 (source0.tex lines 847--857). The article's complete original proof of it is delivered with it (source0.tex lines 859--909, one \texttt{proof} environment of 51 lines). No mathematics is written, repaired or re-derived: the statement and the proof are the article's own lines, in the article's own notation, and no retained line is substituted or re-flowed.  Retained uncounted as the source-local context the proof needs: lines 177--179; lines 749--755; lines 765--767; lines 787--789; lines 817--825.  Nothing was omitted from any retained span, so the package carries no omission record.  No retained line contains a citation, so the wrapper carries no bibliography and no bibliography entry.  No retained line contains a figure environment, an image-inclusion command or drawing code, so no image is delivered and the package declares no asset dependency.  Editorial formatting only, in addition: an AMS \texttt{amsart} preamble that restates the article's own declarations and macros used by the retained lines, namely the environments \texttt{corollary}, \texttt{lemma}, \texttt{proposition} on separate counters, together with the standard packages those lines need.  The retained environments print in this document as Corollary 1, Lemma 1, Proposition 1 in order of appearance, whereas the article prints the same items with the numbering of its own full document.  The complete native source of the article as distributed in the arXiv e-print for this exact version is bundled unchanged as \texttt{sources/196162/source0.tex}.
	\begin{corollary}\label{cor:1}
		Every shedding vertex of $\Delta$ is dismissing.  If $\Delta$ is pure, then a vertex is dismissing if and only if it is shedding.
	\end{corollary}

	\begin{corollary}\label{cor:7}
		If $G=G_1\mathbin{\dot\cup}\cdots\mathbin{\dot\cup}G_s$, then
		\[
		i(G)=\sum_{j=1}^s i(G_j),
		\]
		and $\Ind(G)$ is strongly vertex dismissible if and only if every $\Ind(G_j)$ is strongly vertex dismissible.
	\end{corollary}

	\begin{lemma}\label{lem:5}
		If $\mathcal T$ is a forest, then $\Ind(\mathcal T)$ is vertex decomposable and hence strongly vertex dismissible.
	\end{lemma}

	\begin{lemma}\label{lem:6}
		For every integer $q\geq1$ and every field $\K$, the pure complex $\Ind(C_{3q+1})^{[q]}$ is not Cohen--Macaulay over $\K$.
	\end{lemma}

	\begin{lemma}\label{lem:7}
		For every rooted graph $(H,r)$, we have $i(H)=\min\{1+\alpha(H,r),\beta(H,r)\}$, $i(H-r)=\min\{\beta(H,r),\gamma(H,r)\}$, and $\alpha(H,r)\leq \gamma(H,r)$.  Therefore,
		\[
		r\text{ is not dismissing in }\Ind(H)
		\quad\Longleftrightarrow\quad
		\gamma(H,r)=\alpha(H,r)<\beta(H,r).
		\]
		If two rooted graphs $H_1$ and $H_2$ meet only at their common root $r$, then $r$ is not dismissing in $\Ind(H_1\cup H_2)$ if and only if it is not dismissing in both $\Ind(H_1)$ and $\Ind(H_2)$.
	\end{lemma}

	\begin{proposition}\label{prop:7}
		Let $G$ be a connected unicyclic graph with unique cycle $C=c_0c_1\cdots c_{m-1}c_0$. After deleting the cycle edges, let $(T_j,c_j)$ be the rooted tree component containing $c_j$, and put $\Delta=\Ind(G)$.  The following are equivalent:
		\begin{enumerate}[label=\textup{(\roman*)}]
			\item either $m\not\equiv1\pmod3$, or $m=3q+1$ and some $c_j$ is dismissing in $\Ind(T_j)$;
			\item $\Delta$ is strongly vertex dismissible;
			\item $\Delta$ is vertex dismissible;
			\item $\Delta$ is scalable;
			\item $\Delta$ is initially Cohen--Macaulay over some field;
			\item $\Delta$ is initially Cohen--Macaulay over every field.
		\end{enumerate}
	\end{proposition}

	\begin{proof}
		We first prove (i) implies (ii).  Suppose that $m\not\equiv1\pmod3$ and induct on $|V(G)|$.  If $G=C_m$, then for every $v\in V(C_m)$,
		\[
		\del_{\Ind(C_m)}(v)=\Ind(P_{m-1}),
		\qquad
		\link_{\Ind(C_m)}(v)=\Ind(P_{m-3}),
		\]
		and
		\[
		i(C_m-v)=\left\lceil\frac{m-1}{3}\right\rceil
		\geq \left\lceil\frac m3\right\rceil=i(C_m)
		\]
		exactly when $m\equiv0,2\pmod3$.  Thus $v$ is dismissing, and both branches are strongly vertex dismissible by Lemma~\ref{lem:5}.
		
		If $G$ is not a cycle, choose a leaf $u$ with neighbor $v$.  The vertex $v$ is shedding in $\Ind(G)$ and therefore dismissing by Corollary~\ref{cor:1}.  Every component of $G-v$ and $G-N_G[v]$ is a forest or a smaller connected unicyclic graph with the same cycle.  The induction hypothesis, Lemma~\ref{lem:5}, and Corollary~\ref{cor:7} show that both resulting independence complexes are strongly vertex dismissible.
		
		Now let $m=3q+1$ and suppose that some $c_j$ is dismissing in $\Ind(T_j)$.  View $G$ as the rooted one-sum of $T_j$ and the rest of $G$, both rooted at $c_j$.  By Lemma~\ref{lem:7}, the vertex $c_j$ is dismissing in $\Ind(G)$.  Both $G-c_j$ and $G-N_G[c_j]$ are forests.  Apply Lemma~\ref{lem:5} and Corollary~\ref{cor:7}.
		
		The implications already proved give
		\[
		\textup{(ii)}\Longrightarrow\textup{(iii)}
		\Longrightarrow\textup{(iv)}
		\Longrightarrow\textup{(vi)}
		\Longrightarrow\textup{(v)}.
		\]
		Assume that (i) fails.  Then $m=3q+1$ and every $c_j$ is not dismissing in $\Ind(T_j)$.  Write
		\[
		\alpha_j=\alpha(T_j,c_j),\qquad
		\beta_j=\beta(T_j,c_j),\qquad
		\gamma_j=\gamma(T_j,c_j).
		\]
		By Lemma~\ref{lem:7}, $\gamma_j=\alpha_j<\beta_j$.  Put $k_j=\alpha_j$ and $\kappa=\sum_jk_j$.  We claim that
		\begin{equation}\label{eq:1}
			i(G)=\kappa+q+1.
		\end{equation}
		A selected root contributes at least $k_j+1$, an excluded root dominated inside its tree contributes at least $\beta_j\geq k_j+1$, and a root undominated inside its tree contributes at least $\gamma_j=k_j$ and must have a selected cycle neighbor.  If $\mathcal S$ is the set of selected roots and $\mathcal B$ is the set of roots dominated internally, then $\mathcal S$ is independent and every root outside $\mathcal S\cup\mathcal B$ has a neighbor in $\mathcal S$.  Therefore
		\[
		3q+1\leq3|\mathcal S|+|\mathcal B|,
		\]
		so $|\mathcal S|+|\mathcal B|\geq q+1$.  This proves the lower bound in~\eqref{eq:1}. For the reverse inequality, combine an independent dominating set of $C_{3q+1}$ of size $q+1$ with minimum sets of state $\alpha$ at the selected roots and minimum sets of state $\gamma$ at the remaining roots.
		
		Choose minimum sets $D_j$ of state $\gamma$ and put $D=\bigcup_jD_j$.  Then $D$ is independent, $|D|=\kappa$, it dominates every vertex outside the cycle, and it dominates no cycle vertex.  Hence
		\[
		G-N_G[D]=C_{3q+1}.
		\]
		Let $\Gamma=\Delta^{[i(G)-1]}$.  Since $|D|=\kappa$ and $i(G)=\kappa+q+1$, a facet of $\link_\Gamma(D)$ is obtained exactly by adjoining to $D$ an independent $(q+1)$-set of $G-N_G[D]=C_{3q+1}$.  Conversely, every independent $(q+1)$-set of this remaining cycle gives an independent set of $G$ of size $i(G)$ and hence a facet of $\Gamma$ containing $D$.  Therefore
		\[
		\link_\Gamma(D)=\Ind(C_{3q+1})^{[q]}.
		\]
		By Lemma~\ref{lem:6}, this link is not Cohen--Macaulay over $\K$ for every field $\K$.  Hence $\Gamma$ is not Cohen--Macaulay, contradicting (v).  Thus (v) implies (i).
	\end{proof}

\end{document}
